\chapter*{Resumo}
\addcontentsline{toc}{chapter}{Resumo}

A crescente adoção de AMRs nas indústrias de logística, intralogística,
retalho e manufatura é uma motivação para melhorar a robustez dos
algoritmos de Simultaneous Localization and Mapping (SLAM) em condições
reais de operação, nas quais o ambiente se encontra continuamente
afetado por pessoas, veículos e equipamentos em movimento. As formulações
clássicas de SLAM~2D, fundamentadas no pressuposto de um mundo estático,
tendem a produzir mapas distorcidos ou esbatidos e estimativas instáveis
da pose do robô quando estas observações dinâmicas são integradas sem as
devidas considerações. Esta Tese de Doutoramento, intitulada
\textit{RicoSLAM: localization and mapping in dynamic environments},
investiga como a dinâmica do ambiente pode ser integrada nos processos
de localização e mapeamento de um sistema de SLAM~2D, recorrendo
exclusivamente a odometria das rodas e a um \emph{scanner} laser 2D como
modalidades de perceção.

A investigação organiza-se em torno de três contribuições principais.
Primeiro, propõe-se um detetor de pontos de interesse do tipo canto para
dados de \emph{scanner} laser 2D, em que o ponto-chave resulta da
interseção em forma fechada de duas linhas ajustadas localmente em
múltiplas escalas de vizinhança baseadas em distância. Os resultados
experimentais evidenciam a invariância ao ponto de vista e à geometria
do ambiente das características extraídas. Segundo, introduz-se uma
estrutura de registo de nuvens de pontos 2D baseada em mapas de
distância, que substitui a pesquisa explícita do vizinho mais próximo
característica do Iterative Closest Point (ICP) clássico por um acesso
até $O(1)$ numa matriz de distâncias pré-calculada.
Derivam-se formulações ponto-a-ponto e ponto-a-plano sobre o \emph{manifold}
$SE(2)$, no contexto do método de Gauss-Newton, recorrendo a uma
perturbação à direita, com suporte para múltiplos níveis de
representação e organizações de armazenamento esparsas ou densas.
Por fim, desenvolve-se a framework RicoSLAM propriamente dita, que
integra o \emph{tracker} baseado em mapas de distância como
\emph{front-end} com um \emph{back-end} de grafo de fatores suportado
pela biblioteca srrg2\_solver~\cite{grisetti:robotics:2020}.
A framework inclui critérios de mudança e integração de dados baseados na razão
de \emph{inliers} para a gestão de \emph{keyframes}, um mecanismo de deteção
de visitas a lugares já mapeados combina proximidade no referencial global com
distância em grafo, e uma política de relocalização que
transfere o \emph{tracker} para um \emph{keyframe} existente do grafo de
forma a limitar o seu crescimento. Dois mecanismos complementares
abordam a dinâmica do ambiente: uma fusão \emph{online} baseada em \emph{ray tracing}
que mantém o mapa local de cada \emph{keyframe} filtrando retornos \emph{see-through}
e adicionando novas observações, e uma etapa \emph{offline} baseada em \emph{ray tracing}
para suprimir observações dinâmicas durante a geração da grelha de
ocupação~2D. Adicionalmente, propõem-se fatores em otimização por grafo de
Global Bundle Adjustment (GBA) baseados em mapas de distância, que
complementam as restrições de pose-pose do grafo \emph{online} com
fatores por ponto, derivados nas variantes ponto-a-ponto e
ponto-a-plano para alinhamento de núvens de pontos 2D.

As contribuições propostas são avaliadas no \emph{dataset}
IILABS~3D~\cite{ribeiro:access:2025} e comparadas com os algoritmos
GMapping~\cite{grisetti:tro:2007} e
SLAM~Toolbox~\cite{macenski-jambrecic:joss:2021}. O RicoSLAM atinge o
menor Absolute Trajectory Error (ATE) nas quatro sequências planares
consideradas, com a correção de pose \emph{offline} a reduzir ainda mais
o erro nas sequências com múltiplas revisitas ao mesmo local.
A ativação da política de relocalização mantém o grafo do
\emph{back-end} compacto, com uma dimensão que depende essencialmente da
extensão espacial do ambiente explorado e não do tempo de operação ou
da distância percorrida. O mecanismo \emph{offline} de remoção da
dinâmica é avaliado em sequências adicionais de datasets de AMRs em
ambientes industriais e académico (Avantgarde, Cappero, Flowbotic e Kuka),
melhorando consistentemente a distribuição empírica de ocupação da grelha~2D
resultante e reclassificando para um estado \emph{desconhecido} as
células cujas únicas observações de suporte correspondiam a observações dinâmicas.
Por fim, um \emph{benchmark} sintético valida os fatores de GBA propostos como
uma etapa de otimização global do grafo de poses do \emph{back-end}.

\hfill

\textbf{Palavras-chave:}
ambientes dinâmicos;
deteção de pontos de interesse;
\emph{global bundle adjustment};
localização e mapeamento simultâneos;
mapas de distância;
otimização de grafo de poses;
registo de nuvens de pontos;
robôs móveis autónomos;
\emph{scanner} laser 2D.

\cleardoublepage

\chapter*{Abstract}
\addcontentsline{toc}{chapter}{Abstract}

The growing deployment of AMRs in logistics, intralogistics, retail,
and manufacturing industries is a motivation to improve the robustness of
Simultaneous Localization and Mapping (SLAM) algorithms under realistic
operating conditions, in which the scene is continuously affected by
moving people, vehicles, and equipment. Classical 2D~SLAM formulations,
which rely on a static-world assumption, tend to produce blurred or
distorted maps and unstable robot pose estimates when these dynamic
returns are integrated without proper considerations. This PhD~Thesis,
entitled
\textit{RicoSLAM: localization and mapping in dynamic environments},
investigates how environment dynamics can be incorporated into both
localization and mapping in a 2D~SLAM framework using wheeled
odometry and 2D~laser scanner data as the only sensing modalities.

The research is organized around three main contributions.
First, a line-fitting-based corner-like feature detector for 2D~laser
scanner data is proposed, in which the keypoint is computed as the
closed-form intersection of two locally fitted lines at multiple
distance-based neighborhood scales. Experimental results show that the proposed
features exhibit viewpoint and geometric invariance suitable for
feature-based tracking modules. Second, a 2D~point cloud registration
framework based on unsigned distance maps is introduced, replacing the
explicit nearest-neighbor search of classical
Iterative Closest Point (ICP) with an up to $O(1)$ lookup on a precomputed
distance field. Point-to-point and point-to-plane error formulations are derived
within the Gauss-Newton framework on the $SE(2)$ manifold via right-sided
perturbation, supporting multiple representation levels and sparse or
dense storage layouts. Third, the RicoSLAM framework itself is
developed, integrating the distance map-based tracker as the front-end
with a factor-graph back-end built on top of the
srrg2\_solver~\cite{grisetti:robotics:2020} library. The framework
includes inlier-ratio-based split and merge triggering criteria for
keyframe management, a loop closure detection procedure that combines
world-frame proximity with a graph-distance constraint, and a
relocalization policy that transfers the tracker into an existing
keyframe of the back-end graph to bound its growth. Two complementary
mechanisms address scene dynamics: an online ray-based merger that
maintains the local map of each keyframe by filtering see-through and
add-new returns, and an offline ray-based stage to suppress dynamic
observations during 2D~occupancy grid map generation. In addition,
distance map-based Global Bundle Adjustment (GBA) factors are
proposed, augmenting the pose-pose constraints of the online graph with
per-point distance map factors derived in both point-to-point and
point-to-plane variants.

The proposed contributions are evaluated on the
IILABS~3D~\cite{ribeiro:access:2025} dataset against the
GMapping~\cite{grisetti:tro:2007} and
SLAM~Toolbox~\cite{macenski-jambrecic:joss:2021} baselines. RicoSLAM
attains the lowest Absolute Trajectory Error (ATE) across the four
planar sequences considered, with an offline pose-correction stream that further
reduces the error on sequences with multiple revisits on the same location. The
activation of the relocalization policy keeps the back-end graph
compact, with a size that depends primarily on the spatial extent of
the explored environment rather than on the operation time or path length.
The offline dynamics removal mechanism is additionally evaluated on industrial
AMR deployment sequences from the Avantgarde, Cappero, Flowbotic,
and Kuka datasets, consistently sharpening the empirical occupancy
distribution of the resulting 2D~grid map and reclassifying to an
\emph{unknown} state cells whose only supporting observations were transient
dynamic returns. Lastly, a synthetic Manhattan-like benchmark validates
the proposed distance map-based GBA factors as a globally
consistent refinement stage for the back-end pose graph.

\hfill

\textbf{Keywords:}
2D~laser scanner;
autonomous mobile robots;
distance maps;
dynamic environments;
feature detection;
global bundle adjustment;
point cloud registration;
pose graph optimization;
simultaneous localization and mapping.

\cleardoublepage
